\chapter{Density, Ahlfors bounds and the canonical representative}
\label{ch:density}
\module{NoCompromise.Regularity.Density,
        NoCompromise.Regularity.DensityAhlfors,
        NoCompromise.Regularity.DensityCutComparison,
        NoCompromise.Regularity.DensityEstimates,
        NoCompromise.Regularity.DensityRadial,
        NoCompromise.Regularity.RepresentativeBoundary,
        NoCompromise.Regularity.RepresentativeBounded}

\begin{notation}\label{not:mx-nx}
\uses{def:omega-minimal,def:density}
Throughout this chapter $E$ is $\omega$-minimal
(Definition~\ref{def:omega-minimal}).  For $x\in\ebd E$ put
\[
  m_x(r):=|E\cap B_r(x)|,\qquad n_x(r):=|B_r(x)\setminus E| .
\]
\end{notation}

\begin{lemma}[Coarea derivatives]\label{lem:mx-derivative}
\uses{not:mx-nx,prop:cut-identities}
For a.e.\ $r>0$,
\begin{equation}\label{eq:mx-derivative}
  m_x'(r)=\Hs^2\bigl(E^{(1)}\cap\partial B_r(x)\bigr),
  \qquad
  n_x'(r)=4\pi r^2-m_x'(r),
\end{equation}
and the cut identities \eqref{eq:cut-in}--\eqref{eq:cut-out} hold at $r$.
\end{lemma}

\begin{proof}\uses{cor:slicing,prop:cut-identities,lem:good-radius-ae}
Corollary~\ref{cor:slicing} and Lemma~\ref{lem:good-radius-ae}.
\end{proof}

\begin{lemma}[Cut comparison]\label{lem:cut-comparison}
\uses{def:omega-minimal,not:mx-nx}
For $x\in\ebd E$ and a.e.\ $0<r\le1$,
\begin{equation}\label{eq:cut-comparison}
  \Per\bigl(E;B_r(x)\bigr)\le m_x'(r)+\omega\,m_x(r),
\end{equation}
and symmetrically, comparing with $E\cup B_r(x)$,
\begin{equation}\label{eq:fill-comparison}
  \Per\bigl(E;B_r(x)\bigr)\le n_x'(r)+\omega\,n_x(r).
\end{equation}
\end{lemma}

\begin{proof}\uses{def:omega-minimal,lem:mx-derivative,lem:lsc,
                    prop:cut-identities}
Apply \eqref{eq:omega-minimal} to the competitor
$F=E\setminus B_{r-\eta}(x)$, whose symmetric difference from $E$ is compactly
contained in $B_r(x)$, and then let $\eta\downarrow0$.  The BV trace formula
\eqref{eq:cut-out} and Lemma~\ref{lem:lsc} give the same inequality as the
formal comparison with $E\setminus B_r(x)$, namely
\eqref{eq:cut-comparison}.  For \eqref{eq:fill-comparison} use
$F=E\cup B_{r-\eta}(x)$.
\end{proof}

\begin{constants}\label{const:rd}
\uses{not:mx-nx}
In Lemma~\ref{lem:two-sided-density} below, set
\[
 r_d:=
 \begin{cases}
  1,&\omega=0,\\
  \min\{1,c_I/(2\omega|B_1|^{1/3})\},&\omega>0,
 \end{cases}
 \qquad c_I=(36\pi)^{1/3}.
\]
Thus $0<r_d\le1$ and
$\omega|B_1|^{1/3}r_d\le c_I/2$; then $c_d$ depends only on $c_I$.
The order of choice is
$\omega\rightsquigarrow r_d\rightsquigarrow c_d$.
\end{constants}

\begin{lemma}[Two-sided density]\label{lem:two-sided-density}
\uses{def:omega-minimal,not:mx-nx}
There are $r_d=r_d(\omega)>0$ and $c_d=c_d(\omega)\in(0,|B_1|/2)$ such that
for every $x\in\ebd E$ and $0<r<r_d$,
\begin{equation}\label{eq:two-sided-density}
  c_dr^3\le|E\cap B_r(x)|\le\bigl(|B_1|-c_d\bigr)r^3 .
\end{equation}
\end{lemma}

\begin{proof}
\uses{lem:cut-comparison,thm:sharp-isoperimetric,prop:cut-identities,
      lem:mx-derivative,const:rd}
Apply \eqref{eq:sharp-isoperimetric} to $E\cap B_r(x)$ and use
\eqref{eq:cut-in}:
\[
  c_Im_x(r)^{2/3}\le\Per\bigl(E\cap B_r(x)\bigr)
  =\Per\bigl(E;B_r(x)\bigr)+m_x'(r) .
\]
With \eqref{eq:cut-comparison},
\begin{equation}\label{eq:density-ode}
  c_Im_x(r)^{2/3}\le2m_x'(r)+\omega\,m_x(r).
\end{equation}
Since $m_x(r)\le|B_1|r^3$, choosing $r_d$ as in
Constants~\ref{const:rd} absorbs the last term:
\begin{equation}\label{eq:density-ode-clean}
  m_x'(r)\ge c\,m_x(r)^{2/3}\qquad\text{for a.e.\ }r<r_d .
\end{equation}
At an essential-boundary point $m_x(r)>0$ for every $r>0$, so
$\tfrac{d}{dr}m_x(r)^{1/3}=m_x'(r)/(3m_x(r)^{2/3})\ge c/3$ a.e., and
integrating from zero gives $m_x(r)\ge c_dr^3$.  For the complement, use
\eqref{eq:fill-comparison}, apply \eqref{eq:sharp-isoperimetric} to
$B_r(x)\setminus E$, and repeat with $n_x$; then
$n_x(r)\ge c_dr^3$ is the upper bound in \eqref{eq:two-sided-density}.
\end{proof}

\begin{lemma}[Ahlfors perimeter bounds]\label{lem:ahlfors}
\uses{lem:two-sided-density}
After reducing $r_d$ if necessary, there are
$0<c_P<C_P<\infty$, depending only on $\omega$, such that for every
$x\in\ebd E$ and $0<r<r_d/2$,
\begin{equation}\label{eq:ahlfors}
  c_Pr^2\le\Per\bigl(E;B_r(x)\bigr)\le C_Pr^2 .
\end{equation}
\end{lemma}

\begin{proof}
\uses{prop:rel-iso-ball,lem:two-sided-density,lem:cut-comparison,
      lem:mx-derivative}
\emph{Lower bound.}  \eqref{eq:rel-iso-ball} together with
\eqref{eq:two-sided-density}.

\emph{Upper bound.}  Fix $r<r_d/2$.  Averaging \eqref{eq:mx-derivative} over
$(r,2r)$, choose a radius $\rho\in(r,2r)$ at which
\eqref{eq:cut-comparison} holds and
\[
  m_x'(\rho)\le\frac{m_x(2r)-m_x(r)}{r}\le Cr^2 .
\]
Then \eqref{eq:cut-comparison} at $\rho$ gives
$\Per(E;B_\rho(x))\le Cr^2+C\omega r^3\le C_Pr^2$, and monotonicity of the
perimeter measure under inclusion gives the same bound on $B_r(x)$.
\end{proof}

\section{The canonical representative}

\begin{convention}[The representative]\label{conv:representative}
\uses{def:density,def:omega-minimal}
For an $\omega$-minimal set $E$ put
\begin{equation}\label{eq:representative}
  \Omega:=E^{(1)} .
\end{equation}
From this point on, $\Omega$ always denotes the density-one representative of
whichever $\omega$-minimal set is under discussion.  Statements about
$\partial\Omega$ are statements about a genuine topological boundary, not
about a null-set equivalence class.
\end{convention}

\begin{lemma}[Openness and representative conversion]\label{lem:open-representative}
\uses{conv:representative,lem:two-sided-density,lem:ahlfors}
The set $\Omega=E^{(1)}$ is open, $\ind\Omega=\ind E$ a.e., $E^{(0)}$ is open,
and
\begin{equation}\label{eq:boundary-identification}
  \partial\Omega=\ebd E=\overline{\rbd E} .
\end{equation}
Consequently every point of $\partial\Omega$ satisfies
\eqref{eq:two-sided-density} and \eqref{eq:ahlfors}.
\end{lemma}

\begin{proof}
\uses{lem:two-sided-density,lem:ahlfors,prop:rel-iso-ball,
      thm:degiorgi-structure,def:density}
Lebesgue's density theorem gives $\ind\Omega=\ind E$ a.e.  Suppose
$x\in E^{(1)}$ were not interior to $E^{(1)}$.  Then every sufficiently small
ball about $x$ contains positive measure of both $E$ and its complement, so
\eqref{eq:rel-iso-ball} gives positive perimeter in that ball.  Since
$\mu_E=\Hs^2\llcorner\rbd E$ (Theorem~\ref{thm:degiorgi-structure}), the ball
contains a reduced-boundary point, and such points $x_j\to x$ may be chosen.
Put $d_j:=|x_j-x|$; the complement-density estimate in
$B_{d_j/2}(x_j)\subset B_{2d_j}(x)$ gives
$|B_{2d_j}(x)\setminus E|\ge cd_j^3$, contradicting density one at $x$.  So
$E^{(1)}$ is open, and the same argument with the phases exchanged shows
$E^{(0)}$ is open.  Their common complement is therefore both
$\partial\Omega$ and, by definition, $\ebd E$.  Finally
$\ebd E=\overline{\rbd E}$: the lower perimeter-density bound
\eqref{eq:ahlfors} makes every ball centred at an essential-boundary point
have positive perimeter and hence meet $\rbd E$, while the reverse inclusion
is automatic.
\end{proof}

\begin{lemma}[Boundedness]\label{lem:bounded-representative}
\uses{lem:open-representative,lem:two-sided-density}
If $E$ is $\omega$-minimal with $|E|<\infty$ then $\Omega=E^{(1)}$ is bounded,
so $\partial\Omega$ is compact.
\end{lemma}

\begin{proof}\uses{lem:two-sided-density,lem:open-representative}
Suppose not.  Choose $y_j\in\Omega$ tending to infinity and mutually separated
by more than $4r_d$.  The balls $B_{r_d}(y_j)$ are pairwise disjoint after a
harmless thinning, so finite volume implies infinitely many are not contained
in $\Omega$; keep only those indices.  Since $y_j\in\Omega$ and the ball is
convex, a segment from $y_j$ to a point outside $\Omega$ meets
$\partial\Omega$; choose $x_j\in\partial\Omega\cap B_{r_d}(y_j)$.  Thin once
more so that the balls $B_{r_d/4}(x_j)$ are pairwise disjoint.  By
the lower phase-density bound \eqref{eq:two-sided-density},
\[
  |E|\ge\sum_j|E\cap B_{r_d/4}(x_j)|
  \ge\sum_jc_d(r_d/4)^3=\infty,
\]
a contradiction.
\end{proof}

